\section{The finite eigenclass phase criterion}\label{the-finite-eigenclass-phase-criterion}

This complete finite proof from the supplied continuation is the source of NE13--16. The projection in the necessity construction is specified explicitly; NE12--14 give its relation to the fixed programme kernel.

\subsection{7. The supplied mixed phase has an exact eigenclass criterion}\label{the-supplied-mixed-phase-has-an-exact-eigenclass-criterion}

Now use precisely the mixed pair C1 in the new attachment:

\[
z=\langle Pv,A(I-P)v\rangle_G,
\qquad
w=\langle(I-P)v,APv\rangle_G,
\]

where \(P\) is the orthogonal projector onto the \textbf{specified} kernel. The attachment retains both of these pairings and their complete rational numerators.

Assume \(Av=\lambda v\), and put

\[
E=\|v\|_G^2,\qquad
x=v/\sqrt E,\qquad
r=(A^{\dagger_G}-\bar\lambda I)x,
\qquad
\theta=\|Px\|_G^2.
\]

Then \(r\perp x\), and direct expansion proves

\[
w-z=E\langle r,Px\rangle_G.
\]

Therefore

\[
\boxed{
\Re(\bar zw)
=
\left|\frac{z+w}{2}\right|^2
-\frac{E^2}{4}|\langle r,Px\rangle_G|^2.
}
\tag{38}
\]

Since

\[
\|Px-\theta x\|^2=\theta(1-\theta),
\]

the negative part has the sharp bound

\[
\boxed{
[-\Re(\bar zw)]_+
\le
\frac{E^2}{4}\|r\|_G^2\theta(1-\theta)
\le
\frac{E^2}{16}\|r\|_G^2.
}
\tag{39}
\]

This controls the full complex phase defect by the actual adjoint-eigenvector residual, not by a determinant derivative.

\subsubsection{A projection-independent sign is equivalent to a reducing eigenline}\label{a-projection-independent-sign-is-equivalent-to-a-reducing-eigenline}

For every prescribed rank \(1\le m\le q-1\),

\[
\boxed{
\Re(\bar zw)\ge0
\text{ for every rank-}m\text{ orthogonal projection }P
\iff
A^{\dagger_G}v=\bar\lambda v.
}
\tag{40}
\]

The reverse implication is immediate from \(w=z\).

For necessity, suppose \(r\ne0\), write \(\epsilon=\|r\|\), and take

\[
u_0=r/\epsilon,\qquad g=\lambda-\langle u_0,Au_0\rangle.
\]

Rotate the phase of \(u_0\) so that \(\langle x,Au_0\rangle\) points opposite to \(g\), and set

\[
\theta=
\frac12\left(
1-\frac{|g|}{\sqrt{|g|^2+\epsilon^2}}
\right).
\]

Choose \(P\) so that

\[
Px=\theta x+\sqrt{\theta(1-\theta)}\,u_0.
\]

This is a rank-one orthogonal projection on the two-dimensional span; it can be extended to any prescribed rank \(m\) by directions orthogonal to that span. The exact value is

\[
\boxed{
\Re(\bar zw)
=
-\frac{E^2\epsilon^4}{16(|g|^2+\epsilon^2)}<0.
}
\tag{41}
\]

Equation (40) is a statement about \textbf{all} projections of the given rank. It does not say that normality is necessary for positivity at the one fixed original observation.
